%% =====================================================================
%%  Les sons de l'anglais : guide de l'alphabet phonétique (pistes 0.01 et 0.02)
%% =====================================================================
\liminaire{sons}{Les sons de l'anglais}{Lire l'alphabet phonétique en dix minutes}

En anglais, l'orthographe ne dit pas toujours comment un mot se prononce : \en{\textbf{ou}}
ne se lit pas de la même façon dans \en{you}, \en{four} et \en{house}. C'est pourquoi chaque
mot du livre est accompagné de sa prononciation en \textbf{alphabet phonétique international}
(API), entre barres obliques : \en{hello} \ipa{həˈloʊ}. Pas besoin de tout apprendre
aujourd'hui : revenez à ces pages chaque fois qu'un symbole vous arrête.

\begin{paire}[3]
\begin{tcolorbox}[otBase,colback=otNavyTint,colframe=otNavyTint,halign=center]
{\otIPAfont\fontsize{22}{24}\selectfont\color{otNavy}ˈ}\par\vspace{1mm}
{\fontsize{9.6}{12}\selectfont Placé \textbf{avant} la syllabe forte : \en{about} \ipa{əˈbaʊt} se dit
« a-\textbf{BAOUT} ».\par}
\end{tcolorbox}
\begin{tcolorbox}[otBase,colback=otNavyTint,colframe=otNavyTint,halign=center]
{\otIPAfont\fontsize{22}{24}\selectfont\color{otNavy}ˌ}\par\vspace{1mm}
{\fontsize{9.6}{12}\selectfont Une syllabe un peu accentuée, moins que la principale : \en{afternoon} \ipa{ˌæftɚˈnuːn}.\par}
\end{tcolorbox}
\begin{tcolorbox}[otBase,colback=otNavyTint,colframe=otNavyTint,halign=center]
{\otIPAfont\fontsize{22}{24}\selectfont\color{otNavy}ː}\par\vspace{1mm}
{\fontsize{9.6}{12}\selectfont Une voyelle \textbf{longue} : \en{see} \ipa{siː}. Sans ce signe, la voyelle est brève.\par}
\end{tcolorbox}
\end{paire}

\begin{tableausons}{0.01}{Les voyelles et les diphtongues}
\sonsCategorie{Voyelles}
\son{iː}{see}{siː}{« i » long et tendu, lèvres étirées comme pour sourire.}
\son{ɪ}{big}{bɪɡ}{« i » bref et relâché, entre le « i » de « vie » et le « é » d'« été ».}
\son{e}{bed}{bed}{« è » bref, comme dans « fête ».}
\son{æ}{cat}{kæt}{Entre « a » et « è » : dites « è » en ouvrant grand la bouche.}
\son{ʌ}{bus}{bʌs}{« a » bref et sourd, proche du « eu » de « peur ».}
\son{ɑː}{hot}{hɑːt}{« a » long et profond, comme dans « pâte ».}
\son{ɔː}{door}{dɔːr}{« o » ouvert et long, comme dans « porte ».}
\son{ʊ}{good}{ɡʊd}{« ou » bref et relâché.}
\son{uː}{you}{juː}{« ou » long, comme dans « roue ».}
\son{ə}{about}{əˈbaʊt}{« e » très faible, comme dans « le ». Le son le plus fréquent de l'anglais, toujours dans une syllabe non accentuée.}
\son{ɚ}{teacher}{ˈtiːtʃɚ}{« e » faible, langue recourbée vers l'arrière : le \textbf{r} américain des fins de mots.}
\son{ɝː}{bird}{bɝːd}{« eu » prolongé, langue recourbée vers l'arrière.}
\sonsCategorie{Diphtongues : deux voyelles qui glissent l'une vers l'autre}
\son{eɪ}{day}{deɪ}{« é » qui glisse vers « i » : « déï ».}
\son{aɪ}{time}{taɪm}{« aï », comme dans « aïe ».}
\son{ɔɪ}{voice}{vɔɪs}{« oï », comme dans « cow-boy ».}
\son{oʊ}{go}{ɡoʊ}{« o » qui glisse vers « ou » : jamais un « o » pur.}
\son{aʊ}{now}{naʊ}{« aou », comme dans « caoutchouc ».}
\end{tableausons}

\begin{tableausons}{0.02}{Les consonnes qui n'existent pas en français}
\son{θ}{think}{θɪŋk}{Le bout de la langue entre les dents, et on souffle : ni « s » ni « f ».}
\son{ð}{this}{ðɪs}{Même position, mais en faisant vibrer la voix : ni « z » ni « d ».}
\son{h}{hello}{həˈloʊ}{Un souffle léger, toujours audible (voir l'unité 1).}
\son{r}{red}{red}{Langue recourbée vers l'arrière, sans toucher le palais : jamais roulé, jamais raclé.}
\son{ŋ}{sing}{sɪŋ}{Le « ng » de « parking », sans prononcer de « g » à la fin.}
\son{w}{we}{wiː}{« ou » très rapide, comme au début de « oui ».}
\son{j}{yes}{jes}{Le « y » de « yaourt » (attention : ce n'est pas le « j » français).}
\son{ʃ}{she}{ʃiː}{Le « ch » de « chat ».}
\son{ʒ}{usually}{ˈjuːʒuəli}{Le « j » de « jour ».}
\son{tʃ}{chair}{tʃer}{« tch », comme dans « match ».}
\son{dʒ}{job}{dʒɑːb}{« dj », comme dans « jean ».}
\end{tableausons}

\begin{astuce}[Les autres consonnes]
\en{p, b, t, d, k, g, f, v, s, z, m, n, l} se prononcent à peu près comme en français, avec
deux nuances : en début de mot, \en{p}, \en{t} et \en{k} sont suivis d'un petit souffle
(\en{pen}, \en{top}, \en{cat}) ; et en anglais américain, un \en{t} placé entre deux voyelles
se prononce presque comme un « d » rapide : \en{water}, \en{later}.
\end{astuce}
